% ============================================================
%  MOCK INTRODUCTION & BACKGROUND
%  "Quantum Period-Finding Attacks on Symmetric Ciphers:
%   Implementation and Comparative Analysis Using Simon's Algorithm"
% ============================================================

\documentclass[11pt,a4paper]{article}
\usepackage[utf8]{inputenc}
\usepackage{amsmath,amssymb,amsfonts}
\usepackage{hyperref}
\usepackage{cite}
\usepackage[margin=1in]{geometry}

\title{Quantum Period-Finding Attacks on Symmetric Ciphers:\\
Implementation and Comparative Analysis Using Simon's Algorithm}

\author{
  Author One\thanks{Department of Computer Science, University of Example.}
  \and
  Author Two\thanks{Institute for Quantum Information, University of Example.}
}
\date{\today}

\begin{document}
\maketitle

% ============================================================
% SECTION 1: INTRODUCTION
% ============================================================
\section{Introduction}

% --- Broad Context ---
The advent of large-scale quantum computing poses a fundamental challenge
to modern cryptography.  Shor's celebrated algorithm demonstrated that
integer factorisation and discrete logarithm problems---the security
foundations of RSA, Diffie--Hellman, and elliptic-curve
cryptosystems---can be solved in polynomial time on a quantum
computer~\cite{Shor1997}.  This existential threat to public-key
cryptography has driven a decade-long international effort toward
post-quantum standardisation, culminating in NIST's publication of
FIPS~203, 204, and 205 in August~2024~\cite{NIST_PQC2024}.

While public-key schemes face a clear and present danger, the impact of
quantum computation on \emph{symmetric-key} cryptography has long been
considered more modest.  Grover's unstructured search
algorithm~\cite{Grover1996} provides a quadratic speedup for exhaustive
key search, effectively halving the security level of any block cipher.
The conventional wisdom holds that simply doubling the key length---for
instance, migrating from AES-128 to AES-256---suffices to maintain
adequate security margins against a quantum
adversary~\cite{NIST_PQC2024, Bernstein2009}.

% --- Specific Problem ---
This comfortable picture, however, rests on an implicit assumption: that
generic search is the \emph{only} meaningful quantum speedup available
against symmetric constructions.  A series of striking results beginning
with Kuwakado and Morii~\cite{Kuwakado2010, Kuwakado2012} has shown this
assumption to be false.  When an adversary is granted quantum
superposition access to a cryptographic oracle---the so-called Q2
model~\cite{BonehZhandry2013, Gagliardoni2016}---certain classical
symmetric schemes admit \emph{exponential} quantum speedups that far
exceed Grover's quadratic barrier.  The enabling mechanism is not
unstructured search but rather \emph{quantum period finding}, most
notably via Simon's algorithm~\cite{Simon1997}.

% --- Existing Approaches (Literature Review) ---
Simon's algorithm, introduced in 1994 and refined in 1997, was the first
quantum algorithm to demonstrate a provable exponential separation
between quantum and classical query complexity in an oracle
setting~\cite{Simon1997}.  It solves the hidden-period problem---given a
function $f:\{0,1\}^n\to\{0,1\}^n$ with the promise $f(x) = f(x \oplus
s)$ for a secret $s$, find $s$---in $O(n)$ quantum queries, whereas any
classical probabilistic algorithm requires
$\Omega(2^{n/2})$~\cite{Simon1997}.  Originally viewed as a
complexity-theoretic curiosity, Simon's algorithm has since found
powerful applications in symmetric-key cryptanalysis.

Kuwakado and Morii pioneered this line of work by constructing a quantum
polynomial-time distinguisher for 3-round Feistel ciphers~\cite{Kuwakado2010},
then extending it to a full key-recovery attack on the Even--Mansour
construction~\cite{Kuwakado2012}.  The Even--Mansour cipher,
$E_{k_1,k_2}(x) = P(x \oplus k_1) \oplus k_2$ for a public permutation~$P$,
is a foundational ``minimalist'' block cipher~\cite{EvenMansour1997,
Dunkelman2012} whose classical security was considered well understood.
The quantum attack constructs $f(x)= E(x) \oplus P(x)$, which satisfies
Simon's promise with period $s = k_1$, enabling key recovery in $O(n)$
queries.

At CRYPTO~2016, Kaplan, Leurent, Leverrier, and Naya-Plasencia
dramatically expanded the scope of these
attacks~\cite{Kaplan2016}.  They showed that Simon's algorithm breaks a
wide array of classically secure symmetric constructions in the Q2 model,
including CBC-MAC, PMAC, GMAC, GCM, OCB, and numerous CAESAR
candidates.  Their work also introduced the first quantum \emph{slide
attack}, demonstrating an exponential acceleration of the classical slide
technique of Biryukov and Wagner~\cite{Biryukov1999} when applied to
key-alternating ciphers with identical round keys.  Classically, the
slide attack requires $O(2^{n/2})$ known plaintexts to find a ``slide
pair'' (two plaintext--ciphertext pairs offset by one round); the quantum
variant recovers the key in $O(n)$ queries regardless of the number of
rounds.

Subsequent work has extended quantum slide attacks in several directions.
Bonnetain, Naya-Plasencia, and Schrottenloher~\cite{Bonnetain2019}
generalised the approach to \emph{sliding with a twist},
\emph{complementation slide}, and \emph{mirror slidex} techniques,
achieving up to four-round self-similarity with cost at most quadratic in
the block size.  Leander and May~\cite{Leander2017} combined Grover's
search with Simon's period finding to break the FX construction (key
whitening), showing that adding whitening keys does not significantly
improve quantum security.  Dong, Dong, and Wang~\cite{Dong2020}
converted classical advanced slide attacks to their quantum analogues,
breaking 2/4K-Feistel and 2/4K-DES in polynomial time and improving
attacks on the Russian GOST standard.

The Luby--Rackoff theorem~\cite{LubyRackoff1988} establishes that a
3-round Feistel network with pseudorandom round functions is a
pseudorandom permutation against classical chosen-plaintext attacks.
Kuwakado and Morii's result~\cite{Kuwakado2010} shows that this guarantee
completely collapses in the Q2 model, as a polynomial-time quantum
distinguisher exists for exactly 3~rounds.

% --- Limitation / Gap / Tension ---
Despite the theoretical significance of these results, their practical
implications remain incompletely understood.  Bonnetain and
Jaques~\cite{Bonnetain2020} presented the first concrete resource
estimates for quantum period-finding attacks against real ciphers
(Chaskey, PRINCE, Elephant), concluding that the qubit counts are
comparable to those required for breaking RSA-2048---a sobering
observation.  However, few works provide end-to-end implementations of
these attacks as executable quantum circuits, making it difficult for
practitioners to build intuition about their mechanics and for educators
to convey their significance.  Moreover, the comparative structure of
these attacks---how the same algorithmic primitive (Simon's algorithm)
manifests differently across Even--Mansour, Feistel, and slide
settings---has not been presented in a unified, reproducible framework.

% --- Pivot / Motivation for New Approach ---
In this work, we bridge this gap by providing a complete, executable
implementation of quantum period-finding attacks on three canonical
symmetric constructions: the Even--Mansour cipher, the 3-round Feistel
network, and the iterated key-alternating cipher (via quantum slide
attack).  All three attacks are unified under a common framework centred
on Simon's algorithm, implemented as quantum circuits simulated using the
Qiskit statevector backend~\cite{Qiskit2024}.  This approach enables
direct comparison of oracle construction techniques, query complexity, and
key-recovery procedures across fundamentally different cipher
architectures.

% --- Scope and Assumptions ---
We restrict our analysis to the Q2 chosen-plaintext model, where the
adversary can query the encryption oracle in quantum
superposition~\cite{BonehZhandry2013}.  Block sizes are kept small
(4--8~bits) for simulation feasibility on classical hardware;
nonetheless, the algorithmic structure and query complexity of all
attacks scale identically to their asymptotic descriptions.  We do not
consider Grover-based attacks, chosen-ciphertext models, or
fault-tolerant quantum hardware constraints, as these lie outside the
scope of the present study.

% --- Contribution Statement ---
Our contributions are as follows:
\begin{enumerate}
    \item We provide the first unified, open-source implementation of
    Simon's algorithm applied to three distinct symmetric cipher
    families---Even--Mansour, 3-round Feistel, and iterated slide
    ciphers---using a common quantum oracle abstraction.

    \item We demonstrate, via executable simulation, that the quantum
    slide attack on an $r$-round iterated cipher achieves $O(n)$ key
    recovery independently of~$r$, providing concrete evidence that
    increasing round count is not a defence against quantum adversaries
    with access to the round function.

    \item We present a comparative analysis of oracle construction
    techniques (truth-table synthesis vs.\ direct CNOT construction) and
    their circuit-level costs, offering practical guidance for
    quantum cryptanalysis implementations.
\end{enumerate}

% --- Significance and Impact ---
These results have implications for both the theoretical understanding
and practical evaluation of post-quantum symmetric security.  Our
framework serves as a reference implementation for researchers studying
quantum attacks on block ciphers, and as a pedagogical tool for
conveying the power of quantum period finding beyond the standard
Grover-centric narrative that dominates current post-quantum security
guidance.

% --- Roadmap ---
The remainder of this paper is organised as follows.
Section~\ref{sec:background} establishes the mathematical background and
notation.  Section~\ref{sec:methods} describes our implementation
methodology, including oracle construction and simulation infrastructure.
Section~\ref{sec:results} presents experimental results across all three
attack families.  Section~\ref{sec:discussion} discusses implications and
limitations.  Section~\ref{sec:conclusion} concludes.


% ============================================================
% SECTION 2: BACKGROUND
% ============================================================
\section{Background}\label{sec:background}

% --- Define the Physical System / Model ---
\subsection{Quantum Computing Primitives}

We work within the standard circuit model of quantum
computation~\cite{NielsenChuang2010}.  A quantum computer operates on
$n$-qubit registers, each inhabiting the Hilbert space
$\mathcal{H}=(\mathbb{C}^2)^{\otimes n}$, with dimension $2^n$.  A
computation consists of preparing an initial state
$|0\rangle^{\otimes n}$, applying a sequence of unitary gates (drawn from
a universal gate set, e.g., $\{H, \text{CNOT}, T\}$), and performing
projective measurements in the computational basis.

A quantum oracle for a classical function $f:\{0,1\}^n\to\{0,1\}^m$ is a
unitary operator $U_f$ acting on $n+m$ qubits as
\begin{equation}\label{eq:oracle}
    U_f |x\rangle|y\rangle = |x\rangle|y \oplus f(x)\rangle,
\end{equation}
where $\oplus$ denotes bitwise XOR.  In the Q2 attack model, the
adversary can query $U_f$ in quantum superposition---that is, prepare an
arbitrary superposition over input states and observe the
output~\cite{BonehZhandry2013, Gagliardoni2016}.

% --- Establish Notation ---
\subsection{Notation}

We adopt the following conventions throughout.
\begin{itemize}
    \item Bitstrings are written in MSB-first (big-endian) order:
    ``$1010$'' denotes the integer~10.
    \item $\oplus$ denotes bitwise XOR over $\{0,1\}^n$.
    \item $P:\{0,1\}^n\to\{0,1\}^n$ denotes a publicly known
    permutation; $P^{-1}$ its inverse.
    \item $S:\{0,1\}^n\to\{0,1\}^n$ denotes a publicly known substitution
    box (S-box), also a permutation.
    \item $k, k_1, k_2 \in \{0,1\}^n$ denote secret keys.
    \item $H = \frac{1}{\sqrt{2}}\begin{psmallmatrix}1&1\\1&-1
    \end{psmallmatrix}$ is the Hadamard gate.
    \item $H^{\otimes n}$ is the $n$-fold tensor product of~$H$.
    \item $\text{GF}(2)$ is the Galois field with two elements;
    all linear algebra is performed modulo~2.
\end{itemize}

% --- Present Core Equations ---
\subsection{Simon's Problem and Algorithm}

\textbf{Simon's problem.}\quad Given oracle access to
$f:\{0,1\}^n\to\{0,1\}^n$ with the promise that there exists
$s\in\{0,1\}^n$ such that
\begin{equation}\label{eq:simon_promise}
    f(x) = f(y) \iff x \oplus y \in \{0^n, s\},
\end{equation}
determine~$s$.  When $s=0^n$, the function is one-to-one; otherwise it
is two-to-one~\cite{Simon1997}.

\textbf{Simon's circuit.}\quad One round of Simon's algorithm proceeds
as follows:
\begin{equation}\label{eq:simon_circuit}
    |0^n\rangle|0^n\rangle
    \xrightarrow{H^{\otimes n}\otimes I}
    \frac{1}{\sqrt{2^n}}\sum_{x}|x\rangle|0^n\rangle
    \xrightarrow{U_f}
    \frac{1}{\sqrt{2^n}}\sum_{x}|x\rangle|f(x)\rangle
    \xrightarrow{H^{\otimes n}\otimes I}
    \text{(measure first register)}.
\end{equation}
Measurement yields a uniformly random $y\in\{0,1\}^n$ satisfying
$y\cdot s = 0 \pmod{2}$, where $y\cdot s = \bigoplus_{i} y_i \wedge
s_i$ is the inner product over $\text{GF}(2)$~\cite{NielsenChuang2010}.

\textbf{Classical post-processing.}\quad After collecting $O(n)$
linearly independent equations $\{y_i \cdot s = 0\}$, Gaussian
elimination over $\text{GF}(2)$ recovers $s$ from the null space of the
resulting matrix~\cite{Simon1997}.  Note that standard floating-point
linear algebra (e.g., \texttt{numpy.linalg.matrix\_rank}) yields
incorrect results over $\text{GF}(2)$, as $[1,1]=[1,0]+[0,1]\pmod{2}$
but these are linearly independent over~$\mathbb{R}$.

% --- Explain Key Mechanisms / Algorithms ---
\subsection{Target Cipher Constructions}

We consider three canonical symmetric constructions, each broken by a
distinct application of Simon's algorithm.

\subsubsection{The Even--Mansour Cipher}

The Even--Mansour cipher~\cite{EvenMansour1997} is defined as
\begin{equation}\label{eq:em}
    E_{k_1,k_2}(x) = P(x \oplus k_1) \oplus k_2,
\end{equation}
where $P$ is a publicly known permutation.  Its classical security was
established with a tight $T = \Omega(2^n/D)$ bound for $D$ known
plaintexts~\cite{Dunkelman2012}.  The quantum attack~\cite{Kuwakado2012}
defines $f(x) = E(x) \oplus P(x) = P(x \oplus k_1) \oplus k_2 \oplus
P(x)$, which satisfies Simon's promise with period $s = k_1$.
After recovering $k_1$, the second key follows as $k_2 = E(0) \oplus
P(k_1)$.

\subsubsection{The 3-Round Feistel Network}

A 3-round Feistel network with keyed round function $F(x) = S[x \oplus
k]$ encrypts $(L,R)$ via~\cite{LubyRackoff1988}:
\begin{equation}\label{eq:feistel}
    (L_i, R_i) = (R_{i-1},\ L_{i-1} \oplus F(R_{i-1})), \quad
    i=1,2,3.
\end{equation}
The quantum attack~\cite{Kuwakado2010, Kaplan2016} constructs
$f(x) = E_L(x,0) \oplus E_L(x,1)$, where $E_L$ denotes the left half
of the ciphertext.  Algebraic analysis reveals that $f$ satisfies
Simon's promise with period $s = S[k] \oplus S[1 \oplus k]$.
The recovered period $s$ reduces the key space to a small candidate set,
from which the true key is identified via brute force.

\subsubsection{The Iterated Key-Alternating Cipher (Slide Attack)}

The iterated cipher applies the same round function $r$ times:
\begin{equation}\label{eq:slide}
    E_k(x) = F_k^r(x), \quad F_k(x) = P(x \oplus k).
\end{equation}
Classically, the slide attack~\cite{Biryukov1999, Biryukov2000} finds a
``slide pair'' in $O(2^{n/2})$ queries, independent of $r$.  The quantum
slide attack~\cite{Kaplan2016, Bonnetain2019} constructs
$f(x) = F_k(x) \oplus P(x) = P(x\oplus k) \oplus P(x)$,
which has period $s = k$ since $f(x\oplus k) = P(x) \oplus P(x\oplus k)
= f(x)$.  Simon's algorithm recovers $k$ in $O(n)$ queries, making the
number of rounds~$r$ cryptographically irrelevant.

% --- Clarify Assumptions and Constraints ---
\subsection{Threat Model and Assumptions}

All attacks in this work operate under the Q2 chosen-plaintext attack
(qCPA) model~\cite{BonehZhandry2013, Gagliardoni2016}, wherein the
adversary can query the encryption oracle in quantum superposition.
Formally, the adversary prepares $|x\rangle|0\rangle$ and obtains
$|x\rangle|E_k(x)\rangle$ for all $x$ simultaneously via a single
oracle call.

This model is strictly stronger than the Q1 model (classical queries to
a quantum adversary) and is motivated by scenarios where the
cryptographic device processes quantum states---for example, future
quantum-enabled hardware, or fault-injection settings where the cipher's
internal state is accessible in
superposition~\cite{Kaplan2016,Bonnetain2020}.  We note that Q2 attacks
are not currently realisable on existing quantum hardware; they
constitute an asymptotic threat model for evaluating the long-term
post-quantum security of symmetric designs.

We further assume that all public components ($P$, $S$, circuit
structure) are known to the adversary.  Block sizes in our simulations
range from 4 to 8~bits, constrained by the exponential cost of
statevector simulation ($2^{2n}$ amplitudes for an $n$-bit cipher).
The algorithms themselves are dimension-independent: their $O(n)$
query complexity holds for arbitrary~$n$.

% --- Bridge to Methods ---
\subsection{From Theory to Implementation}

Translating the theoretical attacks described above into executable
quantum circuits requires two key engineering decisions:

\begin{enumerate}
    \item \textbf{Oracle construction.}\quad The quantum oracle $U_f$
    must be synthesised as a reversible circuit.  We employ two
    approaches: (a)~\emph{truth-table synthesis}, which enumerates all
    input--output pairs and implements $U_f$ via multi-controlled X gates
    (MCX), and (b)~\emph{direct CNOT construction}, which exploits the
    algebraic structure of Simon's promise to build the oracle with
    $O(n)$ gates.  Truth-table synthesis is general but scales as
    $O(2^n)$ gates; direct construction is efficient but requires
    knowledge of the period~$s$.

    \item \textbf{Simulation backend.}\quad We use Qiskit's statevector
    simulator~\cite{Qiskit2024}, which provides exact amplitudes and
    deterministic sampling.  This avoids shot noise and allows
    precise verification of the algorithm's correctness, at the cost of
    $O(2^{2n})$ memory.
\end{enumerate}

All code is implemented in Python~3.10+ using Qiskit and NumPy, with
GF(2) linear algebra performed by a custom Gaussian elimination routine
to avoid the pitfalls of floating-point rank computation over
$\text{GF}(2)$.

% --- What Is Not Covered ---
\subsection{Scope Limitations}

We do not consider the following topics, which lie outside the scope of
this work:
\begin{itemize}
    \item \textbf{Grover-based attacks} or hybrid Grover--Simon
    constructions~\cite{Leander2017}.
    \item \textbf{Chosen-ciphertext (Q2-CCA)} or
    \textbf{related-key} attack models.
    \item \textbf{Quantum resource estimation} (T-gate counts,
    error-correction overhead, qubit budgets) for fault-tolerant
    implementations~\cite{Bonnetain2020}.
    \item \textbf{Post-quantum secure symmetric constructions} or
    countermeasures against Q2 attacks.
    \item \textbf{Modes of operation} (CBC-MAC, GCM, OCB) beyond
    the underlying block cipher primitives.
\end{itemize}
We restrict to the core block cipher constructions where the connection
between cipher structure and Simon's promise is most transparent.


% ============================================================
% REFERENCES
% ============================================================
\begin{thebibliography}{99}

\bibitem{Shor1997}
P.~W. Shor,
``Polynomial-time algorithms for prime factorization and discrete
logarithms on a quantum computer,''
\textit{SIAM J. Comput.}, vol.~26, no.~5, pp.~1484--1509, 1997.
\texttt{arXiv:quant-ph/9508027}.

\bibitem{NIST_PQC2024}
National Institute of Standards and Technology,
``Post-Quantum Cryptography Standardization,''
FIPS 203, 204, 205, August 2024.
\url{https://csrc.nist.gov/projects/post-quantum-cryptography}.

\bibitem{Grover1996}
L.~K. Grover,
``A fast quantum mechanical algorithm for database search,''
in \textit{Proc.\ 28th ACM Symposium on Theory of Computing (STOC)},
pp.~212--219, 1996.

\bibitem{Bernstein2009}
D.~J. Bernstein,
``Cost analysis of hash collisions: Will quantum computers make
SHARCS obsolete?''
in \textit{Workshop Record of SHARCS}, 2009.

\bibitem{Kuwakado2010}
H.~Kuwakado and M.~Morii,
``Quantum distinguisher between the 3-round Feistel cipher and the
random permutation,''
in \textit{Proc.\ IEEE International Symposium on Information Theory
(ISIT)}, pp.~2682--2685, 2010.

\bibitem{Kuwakado2012}
H.~Kuwakado and M.~Morii,
``Security on the quantum-type Even-Mansour cipher,''
in \textit{Proc.\ International Symposium on Information Theory and
its Applications (ISITA)}, pp.~312--316, 2012.

\bibitem{BonehZhandry2013}
D.~Boneh and M.~Zhandry,
``Quantum-secure message authentication codes,''
in \textit{Advances in Cryptology---EUROCRYPT 2013},
LNCS, vol.~7881, pp.~592--608, Springer, 2013.

\bibitem{Gagliardoni2016}
T.~Gagliardoni, A.~H\"{u}lsing, and C.~Schaffner,
``Semantic security and indistinguishability in the quantum world,''
in \textit{Advances in Cryptology---CRYPTO 2016},
LNCS, vol.~9816, pp.~60--89, Springer, 2016.

\bibitem{Simon1997}
D.~R. Simon,
``On the power of quantum computation,''
\textit{SIAM J. Comput.}, vol.~26, no.~5, pp.~1474--1483, 1997.

\bibitem{Kaplan2016}
M.~Kaplan, G.~Leurent, A.~Leverrier, and M.~Naya-Plasencia,
``Breaking symmetric cryptosystems using quantum period finding,''
in \textit{Advances in Cryptology---CRYPTO 2016},
LNCS, vol.~9815, pp.~207--237, Springer, 2016.
\texttt{arXiv:1602.05973}.

\bibitem{EvenMansour1997}
S.~Even and Y.~Mansour,
``A construction of a cipher from a single pseudorandom permutation,''
\textit{J. Cryptology}, vol.~10, no.~3, pp.~151--162, 1997.

\bibitem{Dunkelman2012}
O.~Dunkelman, N.~Keller, and A.~Shamir,
``Minimalism in cryptography: The Even-Mansour scheme revisited,''
in \textit{Advances in Cryptology---EUROCRYPT 2012},
LNCS, vol.~7237, pp.~336--354, Springer, 2012.

\bibitem{Biryukov1999}
A.~Biryukov and D.~Wagner,
``Slide attacks,''
in \textit{Fast Software Encryption (FSE 1999)},
LNCS, vol.~1636, pp.~245--259, Springer, 1999.

\bibitem{Biryukov2000}
A.~Biryukov and D.~Wagner,
``Advanced slide attacks,''
in \textit{Advances in Cryptology---EUROCRYPT 2000},
LNCS, vol.~1807, pp.~589--606, Springer, 2000.

\bibitem{Bonnetain2019}
X.~Bonnetain, M.~Naya-Plasencia, and A.~Schrottenloher,
``On quantum slide attacks,''
in \textit{Selected Areas in Cryptography---SAC 2019},
LNCS, vol.~11959, pp.~492--519, Springer, 2019.

\bibitem{Leander2017}
G.~Leander and A.~May,
``Grover meets Simon---Quantumly attacking the FX-construction,''
in \textit{Advances in Cryptology---ASIACRYPT 2017},
LNCS, vol.~10625, pp.~161--178, Springer, 2017.

\bibitem{Dong2020}
X.~Dong, B.~Dong, and X.~Wang,
``Quantum attacks on some Feistel block ciphers,''
\textit{Des.\ Codes Cryptogr.}, vol.~88, pp.~1179--1203, 2020.

\bibitem{LubyRackoff1988}
M.~Luby and C.~Rackoff,
``How to construct pseudorandom permutations from pseudorandom
functions,''
\textit{SIAM J. Comput.}, vol.~17, no.~2, pp.~373--386, 1988.

\bibitem{Bonnetain2020}
X.~Bonnetain and S.~Jaques,
``Quantum period finding against symmetric primitives in practice,''
\textit{IACR Trans.\ Cryptogr.\ Hardw.\ Embed.\ Syst.}, vol.~2022,
no.~1, pp.~1--27, 2022.
\texttt{IACR ePrint 2020/1418}.

\bibitem{NielsenChuang2010}
M.~A. Nielsen and I.~L. Chuang,
\textit{Quantum Computation and Quantum Information},
10th Anniversary Edition.
Cambridge University Press, 2010.

\bibitem{Qiskit2024}
Qiskit Contributors,
``Qiskit: An open-source framework for quantum computing,'' 2024.
\url{https://qiskit.org}.

\bibitem{Zhandry2016}
M.~Zhandry,
``A note on quantum-secure PRPs,''
Cryptology ePrint Archive, Report 2016/1076, 2016.

\end{thebibliography}

\end{document}
